\section{One positive functional, two periods}\label{sec:moments}
To put both periods in the same positive matrix, use Carreira's
functional \cite[Definition~24 and Lemma~25]{Carreira}, before restricting
its resolvent to half-integers. Set
\begin{equation}\label{eq:functional}
 g(y)=\frac1{\sinh(\pi y)},\qquad
 w(y)=-\frac{y^4}{6}g'''(y),\qquad
 \nu(f)=\int_0^\infty f(4y^2)w(y)\dd y.
\end{equation}
The series $g(y)=2\sum_{j\geq0}e^{-(2j+1)\pi y}$ shows that $w>0$.
The weight is bounded near zero and decays exponentially times a
polynomial at infinity. All the integrals below converge.

\subsection{The resolvent, including its rational correction}
For every $a>0$, the inherited identity is
\begin{equation}\label{eq:resolvent}
 \int_0^\infty\frac{w(y)}{y^2+a^2}\dd y
 =a^3\sum_{k\geq0}\frac{(-1)^k}{(k+a)^4}-\frac1{2a}.
\end{equation}
We include its proof because the final rational correction is important
in the later prime calculations. Put $f(y)=y^4/(y^2+a^2)$. Three
integrations by parts, using
$f'''(y)=6a^3\operatorname{Im}(a-iy)^{-4}$, give the left side as
$a^3\int_0^\infty\operatorname{Im}(a-iy)^{-4}g(y)\dd y$.
At zero the three boundary terms are $O(y)$; at infinity they vanish
exponentially. Now use
\[
 \operatorname{Im}(a-iy)^{-4}
 =\frac16\int_0^\infty x^3e^{-ax}\sin(xy)\dd x,
 \qquad
 \int_0^\infty\sin(xy)g(y)\dd y=\frac12\tanh(x/2).
\]
The second formula follows from the exponential series for $g$ and
the logarithmic derivative of Euler's product for $\cosh(x/2)$.
The double integral is absolutely convergent, using $|\sin(xy)|\leq xy$
near zero. Expanding $\tanh(x/2)=1+2\sum_{k\geq1}(-1)^ke^{-kx}$
and integrating proves \eqref{eq:resolvent}.

Write $X$ and $Z$ for indeterminates representing $\bet$ and $\zet$.
For odd $l=2j+1$ and for $m\geq1$, define
\begin{align}
 b_l&=\frac{(-1)^jl^3}{2},&
 r_l^o&=-b_l\sum_{v=0}^{j-1}\frac{(-1)^v}{(2v+1)^4}-\frac1{4l},\label{eq:oddcoeff}\\
 c_m&=\frac{7(-1)^{m-1}m^3}{32},&
 r_m^e&=-\frac{(-1)^{m-1}m^3}{4}
       \sum_{v=1}^{m-1}\frac{(-1)^{v-1}}{v^4}-\frac1{8m}.
       \label{eq:evencoeff}
\end{align}
Empty sums are zero. The change of variable in \eqref{eq:functional}
and \eqref{eq:resolvent} give the formal evaluations
\begin{equation}\label{eq:formal-resolvents}
 \nu\bigl((u+l^2)^{-1}\bigr)=b_lX+r_l^o,\qquad
 \nu\bigl((u+4m^2)^{-1}\bigr)=c_mZ+r_m^e.
\end{equation}
They become genuine integral identities at the actual periods.

\subsection{A fixed cofactor normalization}
For $K\geq1$, put
\begin{gather*}
 a_m=4m^2,\qquad J(u)=\prod_{m=1}^K(u+a_m),\qquad
 O(u)=\prod_{j=0}^{K-1}(u+(2j+1)^2),\\
 D(u)=J(u)O(u),\qquad J_m(u)=\frac{J(u)}{u+a_m}.
\end{gather*}
Here $O$ is a polynomial, not an order symbol. Define the symmetric pencil
\begin{equation}\label{eq:pencil}
 M_K(X,Z)=R_K+XB_K+ZC_K,
 \qquad (M_K)_{mn}=\nu(J_mJ_n/D).
\end{equation}
The right side is evaluated formally by \eqref{eq:formal-resolvents}.
At $(X,Z)=(\bet,\zet)$ it is a positive definite Gram matrix: the
functions $J_m/\sqrt D$ are independent and the measure is positive.

The complete coefficients have a particularly useful form. With $l$
running over the odd integers from $1$ to $2K-1$, set
\begin{align}
 e_m&=\frac{J'(-a_m)}{O(-a_m)}
 =-\frac{4^{2K-1}}
 {2m^2\binom{2K-2m}{K-m}\binom{2K+2m}{K+m}},\label{eq:e}\\
 \rho_l&=\frac{J(-l^2)}{O'(-l^2)},\qquad V_{ml}=\frac1{a_m-l^2}.
 \nonumber
\end{align}
Partial fractions give
\begin{equation}\label{eq:coeffmatrices}
 C_K=\diag(e_mc_m),\quad
 B_K=V\diag(\rho_lb_l)V^T,\quad
 R_K=\diag(e_mr_m^e)+V\diag(\rho_lr_l^o)V^T.
\end{equation}
An even pole contributes only to its corresponding diagonal entry;
an odd pole contributes the rank-one residue $\rho_lV_{\cdot l}V_{\cdot l}^T$.
There is no polynomial quotient here, since the numerator degree is
$2K-2$ and the denominator degree is $2K$.

For even $m$,
\begin{equation}\label{eq:cpositive}
 (C_K)_{mm}=
 \frac{7\,2^{4K-8}m}
 {\binom{2K-2m}{K-m}\binom{2K+2m}{K+m}}>0.
\end{equation}
Thus the even-coordinate principal subspace permits an upper rational
substitution for $\zet$ without losing positivity. We will first
determine its arithmetic exactly.
